\chapter{Межсерийная фотометрия и история объекта}

\section{Назначение Stage 3}

Stage 3 начинается после того, как отдельные серии уже прошли Stage 2 и для них существуют опубликованные \texttt{SCIENCE} generations с таблицей \texttt{science\_measurements}.

Главный принцип архитектуры: межсерийная калибровка вычисляется один раз для поля и совместимой фотометрической системы, а не повторно для каждого объекта.

\begin{lstlisting}
FIELD
  series A ---+
  series B ---+--> inter-series calibration --> object histories
  series C ---+
\end{lstlisting}

После этого долговременная кривая любого объекта строится применением уже найденной поправки соответствующей series.

\section{Необязательный выбор объекта}

Для Stage~1 и Stage~2 не требуется задавать target: поле обрабатывается и калибруется независимо от выбора отдельного объекта. Явный \cli{--target-object-id}, \cli{--target-source-id} либо пара \cli{--target-ra-deg}/\cli{--target-dec-deg} в \cli{process-archive-path} нужна только для downstream target product. Поиск ближайшего к центру объекта или использование \cli{ALPHA/DELTA}, \cli{CRVAL/WCS} либо \cli{coord2000} как target не выполняется; без selector ensemble получает пустой target-exclusion set, а существующие reference-star критерии не меняются.

Без selector target export и отметка \texttt{FIELD\_TARGET} не создаются, существующие target CSV/JSON/PNG не удаляются из-за no-target run, а само отсутствие target не делает поле неготовым. Одно откалиброванное поле может обслуживать несколько последующих object selections: полную Stage~2 нельзя повторять отдельно для каждого target. Immutable BASE и текущая validated SCIENCE сохраняются долговременно; временные и superseded products требуют bounded lifecycle, а reduced Stage~1 FITS reclaimable только по действующей policy после durable Stage~2 validation.

\section{Ограничения первого MVP}

Первый вариант Stage 3 намеренно ограничен наиболее проверяемым случаем:

\begin{itemize}
  \item одно поле;
  \item одинаковый \texttt{instrument\_unit\_id};
  \item одинаковый фильтр;
  \item одинаковый catalog namespace;
  \item минимум две series с опубликованной Stage 2 science photometry;
  \item достаточное число общих объектов.
\end{itemize}

Одинаковый catalog namespace на первом этапе является существенным ограничением. В текущей модели \texttt{objects.object\_uid} включает catalog namespace/source identity. Поэтому совпадение \texttt{object\_id} между independently processed series считается прямым межсерийным identity только при совместимой catalog provenance. Cross-catalog matching будет отдельным поздним этапом Stage 3.

\section{Этап 3.0: проверка готовности поля}

Начиная с версии \texttt{0.5.41.0}, доступна read-only команда:

\begin{lstlisting}
python -m grandgousier.cli plan-interseries-calibration `
  --field-id 12 `
  --min-points-per-object 5 `
  --min-common-objects 20
\end{lstlisting}

Вместо \texttt{--field-id} можно использовать точное имя:

\begin{lstlisting}
python -m grandgousier.cli plan-interseries-calibration `
  --field-name "TIC13533330_01" `
  --filter V
\end{lstlisting}

Команда ничего не записывает в БД. Она отвечает на вопрос, какие уже существующие Stage 2 продукты можно использовать для первой межсерийной калибровки.

Начиная с версии \texttt{0.5.55.0} planner может дополнительно проверить persisted Stage~2 provenance, не повторяя фотометрию:

\begin{lstlisting}
python -m grandgousier.cli plan-interseries-calibration `
  --field-id 33 `
  --catalog-name "auto_catalog_master_20180907_GPX_TF17_2_mixed" `
  --required-stage2-recipe stage3_scale_v1 `
  --min-points-per-object 5 `
  --min-common-objects 15
\end{lstlisting}

Проверка идёт по связи опубликованной \texttt{SCIENCE} generation с создавшим её \texttt{pipeline\_run} и его \texttt{recipe\_json}; имя calibration не используется как научная provenance. Для пригодной series выводится \texttt{READY\_STAGE3\_RECIPE}. Доказанное отличие критического параметра даёт \texttt{RECIPE\_MISMATCH}; отсутствие достаточных persisted полей --- \texttt{MISSING\_PROVENANCE}. Последний статус не интерпретируется по догадке. Результат цветового quality gate, в частности \texttt{color\_calibration\_activated=false}, не является recipe mismatch, если цветовая калибровка была запрошена по recipe.

Canonical \texttt{stage3\_scale\_v1} для новых запусков сохраняет \texttt{ensemble\_min\_ref\_stars\_per\_frame=30}. Readiness-contract допускает исторически real-validated значение 20 как совместимое legacy-состояние наряду с canonical 30 и показывает его отдельной compatibility note; новые runs по-прежнему создаются с 30.

Начиная с версии \texttt{0.5.45.2}, для каждой series сначала выбирается самая новая опубликованная \texttt{SCIENCE} generation по времени финализации продукта: \texttt{finalized\_at}, затем \texttt{materialized\_at}, \texttt{pipeline\_run\_id} и \texttt{calibration\_version\_id}. Номер \texttt{generation\_no} не сравнивается между различными \texttt{logical\_name}, поскольку он нумеруется отдельно внутри каждого логического продукта.

Начиная с clean schema \texttt{0.7.0}, Stage~3 читает финальный durable SCIENCE layer
\texttt{science\_photometry} через единый calibration view. Старые
\texttt{science\_photometry\_cache}/\texttt{science\_measurements} paths относятся
только к frozen 0.6 reference DB и не являются частью active 0.7 workflow.

\subsection{Durable SCIENCE storage в 0.7.0}

В clean schema 0.7 финальный SCIENCE product хранится в узкой таблице
\texttt{science\_photometry}. Она keyed by
\texttt{(calibration\_version\_id, base\_measurement\_id)} и содержит только
скаляры, которые действительно могут изменяться calibration stages:
\texttt{flux}, \texttt{flux\_err}, \texttt{mag\_inst}, \texttt{mag\_cal},
\texttt{mag\_err}, \texttt{mag\_err\_random},
\texttt{mag\_err\_systematic}, \texttt{local\_zp\_correction}.
Геометрия aperture, centroid и object/frame identity не дублируются: они
наследуются через immutable BASE measurement.

Таблицы 0.6 \texttt{science\_measurements},
\texttt{science\_photometry\_cache}, \texttt{science\_flag\_override\_sets} и
\texttt{science\_flag\_overrides} не создаются в новой 0.7 БД. Команды их
backfill/verification/retirement также удалены из active CLI. Они относятся
только к frozen 0.6 reference database и не должны применяться к clean rebuild.

Штатный pipeline не имеет режима, который пишет derived full-copy rows в
\texttt{measurements}: canonical \texttt{measurements} является BASE-only.
Stage~3 и global-object/light-curve readers получают финальную фотометрию через
единый calibration view поверх durable \texttt{science\_photometry}.

Для каждой пары series строится пересечение usable object IDs. Пара считается непосредственно пригодной только если одновременно выполняются условия:

\begin{enumerate}
  \item обе series имеют готовую science generation;
  \item совпадает instrument unit;
  \item совпадает filter;
  \item совпадает catalog namespace;
  \item число общих usable objects не меньше \texttt{--min-common-objects}.
\end{enumerate}

Возможные итоговые статусы planner:

\begin{description}
  \item[\texttt{NEED\_MORE\_READY\_SERIES}] меньше двух пригодных science series;
  \item[\texttt{NO\_ELIGIBLE\_SERIES\_OVERLAP}] series есть, но нет допустимой пары;
  \item[\texttt{READY\_FOR\_PAIRWISE\_MVP}] существует связная допустимая пара;
  \item[\texttt{READY\_FOR\_GRAPH\_VALIDATION}] существует связная группа минимум из трёх series.
\end{description}

Planner также выводит connected components графа series и рекомендует reference series. На этом этапе reference series выбирается детерминированно по числу usable objects, затем по числу science points и frames. Это ещё не photometric quality estimator, а только удобная стартовая точка для следующего алгоритма.

JSON-отчёт можно сохранить командой:

\begin{lstlisting}
python -m grandgousier.cli plan-interseries-calibration `
  --field-id 12 `
  --out reports\stage3_field12_plan.json `
  --format json
\end{lstlisting}

\section{Read-only M1--M3 analysis начиная с 0.5.45.0}

Первая вычислительная реализация Stage~3 запускается командой

\begin{lstlisting}
python -m grandgousier.cli analyze-interseries-calibration `
  --field-id 12 `
  --filter V `
  --out-root reports\stage3_field12_V
\end{lstlisting}

Команда ничего не записывает в БД. Она использует только опубликованные Stage~2 \texttt{SCIENCE} generations и создаёт диагностические CSV/JSON, необходимые для проверки M1--M3 до введения persistent Stage~3 product.

\subsection{M1: representative magnitude объекта в series}

Для объекта $j$ в series $s$ из пригодных точек Stage~2 $m_{ijs}$ используется медиана

\[
  \bar m_{js}=\operatorname{median}_i(m_{ijs}).
\]

Внутрисерийный robust scatter определяется как

\[
  \sigma_{js}^{\mathrm{rob}}=1.4826\,\operatorname{MAD}_i(m_{ijs}),
\]

а неопределённость representative magnitude берётся консервативно как максимум ошибки медианы по robust scatter и типичной measurement uncertainty, делённых на $\sqrt{N}$.

Для inter-series reference population используются объекты ближайшей по provenance финальной ENSEMBLE-итерации с \texttt{is\_reference=true}; стандартные object-QC reject flags дают дополнительный veto. Это не меняет Stage~2 и не создаёт новую reference population скрыто.

\subsection{M2: попарное смещение series}

В версии 0.5.45.0 знак зафиксирован однозначно. Для общего reference object $j$ в series $a$ и $b$:

\[
  d_j^{ab}=\bar m_{jb}-\bar m_{ja}.
\]

После iterative median/MAD clipping оценивается

\[
  \Delta_{ab}=\operatorname{median}_j(d_j^{ab}).
\]

С этой конвенцией одновременно выполняются

\[
  z_a-z_b=\Delta_{ab}
\]

и последующее применение поправки

\[
  m^{\mathrm{global}}_{ijs}=m^{\mathrm{science}}_{ijs}+z_s.
\]

Для каждой пары сохраняются число initial/used/rejected calibrators, robust scatter, uncertainty и полный список calibrators с residuals.

\subsection{M3: глобальное графовое решение}

При наличии трёх и более series pairwise offsets становятся рёбрами графа. После фиксации gauge

\[
  z_{\mathrm{ref}}=0
\]

решается weighted least-squares система

\[
  z_a-z_b\simeq\Delta_{ab}.
\]

Для каждого ребра сохраняется closure residual

\[
  r_{ab}=(z_a-z_b)-\Delta_{ab},
\]

а для всего графа --- $\chi^2$, число степеней свободы, reduced $\chi^2$ и максимальный стандартизованный residual. Это позволяет обнаружить несовместимую series или edge до записи межсерийной калибровки в БД.

Если в одном field scope присутствуют несколько независимых совместимых components, команда не выбирает ``самую большую'' молча. Нужно сузить \texttt{--filter}, \texttt{--instrument-unit-id}, \texttt{--catalog-name} либо явно указать \texttt{--reference-series-id}.

\subsection{Диагностические файлы}

\begin{description}
  \item[\texttt{series\_summary.csv}] identity и качество каждой выбранной series;
  \item[\texttt{series\_object\_states.csv}] M1 representative states;
  \item[\texttt{pairwise\_offsets.csv}] M2 offsets и robust diagnostics;
  \item[\texttt{pairwise\_calibrators.csv}] все использованные и rejected calibrators;
  \item[\texttt{graph\_offsets.csv}] M3 additive offsets $z_s$;
  \item[\texttt{graph\_edges.csv}] closure residuals каждого ребра;
  \item[\texttt{interseries\_summary.json}] полная сводка M1--M3 и provenance input generations.
\end{description}

\section{M4: immutable persistent inter-series generation начиная с 0.5.46.0}

После реальной проверки M1--M3 версия \texttt{0.5.46.0} вводит persistent Stage~3 product. Ранняя таблица \texttt{interseries\_calibration\_offsets} остаётся legacy compatibility layer и не используется как production M4 storage. Новая схема состоит из трёх таблиц:

\begin{description}
  \item[\texttt{interseries\_calibration\_generations}] immutable generation-level provenance и качество графа;
  \item[\texttt{interseries\_calibration\_series}] точные input \texttt{SCIENCE} calibration IDs и additive offsets $z_s$;
  \item[\texttt{interseries\_calibration\_edges}] pairwise M2 diagnostics, membership calibrator object IDs и M3 closure residuals.
\end{description}

Перед первым использованием необходимо применить managed migration:

\begin{lstlisting}
python -m grandgousier.cli schema-status
python -m grandgousier.cli migrate-db
\end{lstlisting}

Generation fingerprint является SHA-256 от точных \texttt{series\_id + calibration\_version\_id} входов, field/instrument/filter/catalog scope, reference series и параметров M1--M3. Поэтому повтор той же команды idempotent, а новая Stage~2 SCIENCE generation или изменение threshold создаёт новый M4 generation. Старый generation не мутируется.

Для redundant графа:

\begin{lstlisting}
python -m grandgousier.cli persist-interseries-calibration `
  --field-id 33 `
  --filter R `
  --catalog-name "auto_catalog_master_20180907_GPX_TF17_2_mixed" `
  --reference-series-id 77 `
  --min-points-per-object 5 `
  --min-common-objects 20
\end{lstlisting}

Если connected graph является деревом, то $\mathrm{dof}=0$ и независимого closure test нет. Такой результат получает reliability \texttt{SPARSE\_GRAPH\_NO\_CLOSURE} и не записывается без явного подтверждения \texttt{--allow-sparse-graph}. Это не понижает threshold автоматически; пользователь должен уже иметь отдельное научное обоснование sparse bridge.

\begin{lstlisting}
python -m grandgousier.cli persist-interseries-calibration `
  --field-id 33 `
  --filter V `
  --reference-series-id 76 `
  --min-common-objects 15 `
  --allow-sparse-graph
\end{lstlisting}

Надёжность успешного generation кодируется явно:

\begin{description}
  \item[\texttt{REDUNDANT\_GRAPH\_OK}] граф имеет положительное число степеней свободы и внутренний closure test;
  \item[\texttt{SPARSE\_GRAPH\_NO\_CLOSURE}] граф связан, но является деревом и closure redundancy отсутствует.
\end{description}

Проверка persisted generation:

\begin{lstlisting}
python -m grandgousier.cli verify-interseries-calibration `
  --generation-id 1
\end{lstlisting}

Verifier не выбирает автоматически более новую Stage~2 calibration. Он проверяет fingerprint, exact source SCIENCE IDs, наличие retained \texttt{science\_measurements}, graph equations, $\chi^2$ и reliability уже записанного generation.

Неопределённость \texttt{offset\_err\_mag} является общей zero-point uncertainty соответствующей series. Она коррелирована между всеми points этой series и поэтому M4 не добавляет её независимо к каждой Stage~2 measurement и не переписывает \texttt{mag\_err\_random}/\texttt{mag\_err\_systematic}. Такое применение относится к следующему object-history слою M5.

\section{M5: read-only global object history начиная с 0.5.47.0}

M5 не материализует новую копию всех measurements в БД. Он читает точные \texttt{science\_measurements}, перечисленные в \texttt{interseries\_calibration\_series}, и применяет только соответствующий persisted M4 offset:
\[
  m_{\mathrm{global},i}=m_{\mathrm{science},i}+z_s.
\]
Join является строгим одновременно по \texttt{series\_id} и \texttt{science\_calibration\_version\_id}. Поэтому M5 не может случайно применить offset к более новой Stage~2 generation.

Базовый вызов:

\begin{lstlisting}
python -m grandgousier.cli export-global-object-light-curve `
  --object-id 12345 `
  --field-id 33 `
  --out-root reports\object_12345_global
\end{lstlisting}

Если в одном scope уже накопилось несколько immutable M4 generations, automatic selection намеренно завершается ошибкой. Для воспроизводимого historical/controlled export нужно явно перечислить generation IDs:

\begin{lstlisting}
  --generation-id 1 `
  --generation-id 2
\end{lstlisting}

Для каждой current SCIENCE series M5 проверяет, совпадает ли newest SCIENCE calibration ID с exact source ID выбранного M4 generation. При наличии current object points в unlinked/stale series export по умолчанию прекращается. \texttt{--allow-partial} разрешает вывести только exact-linked часть, но это должно считаться осознанно неполным продуктом.

\subsection{Выходные файлы}

Для \texttt{--out-root reports\\object\_12345\_global} создаются:

\begin{description}
  \item[\texttt{object\_12345\_global.csv}] точные accepted M5 points; без downsampling;
  \item[\texttt{object\_12345\_global\_summary.json}] M4 provenance, coverage, uncertainty semantics и plot policy;
  \item[\texttt{object\_12345\_global\_V.png}] отдельный plot фильтра V;
  \item[\texttt{object\_12345\_global\_V\_plot.csv}] именно те raw/binned points, которые показаны на plot;
  \item[аналогичные R/g/r/i файлы] создаются независимо для каждого встреченного фильтра.
\end{description}

По умолчанию main CSV содержит только science-accepted points. \texttt{--include-rejected} добавляет rejected rows с \texttt{quality\_flags} и \texttt{reject\_reasons}; plot всё равно строится только по accepted points.

Начиная с 0.5.47.1 случай, когда exact M4 linkage найден, но все rows отвергнуты QC, имеет явный статус \texttt{NO\_ACCEPTED\_POINTS\_AFTER\_QC}. Summary содержит \texttt{rejection\_reasons} и \texttt{rejection\_reasons\_by\_series}. Команда печатает эти счётчики и возвращает exit code 2. Это не означает ошибку M4: например, насыщенная или явно нестабильная звезда может корректно иметь 0 accepted points.

Для детальной проверки можно добавить:

\begin{lstlisting}
--include-rejected-debug
\end{lstlisting}

Тогда создаётся отдельный \texttt{<out-root>\_rejections.csv} с \texttt{series\_id}, \texttt{frame\_id}, exact SCIENCE/M4 IDs, magnitude, merged \texttt{quality\_flags} и \texttt{reject\_reasons}. Этот файл диагностический и не меняет основной accepted CSV. Если accepted points нет, plots намеренно не создаются, поскольку plotting policy использует только science-accepted points.

\subsection{Насыщение в M5 начиная с 0.5.47.3}

Object-QC flag \texttt{SATURATED\_LIKE} имеет series-level смысл: объект имел хотя бы один saturated/near-saturated epoch. Он остаётся важным veto при выборе ensemble reference stars и в Object-QC/candidate policy, но начиная с 0.5.47.3 сам по себе \emph{не} отвергает все epochs объекта при M5 global-history export. M5 использует measurement-level flag \texttt{SATURATED\_OR\_NEAR\_SATURATED} и отвергает только конкретный epoch, на котором этот флаг активен.

Таким образом, наличие \texttt{SATURATED\_LIKE} в \texttt{quality\_flags} accepted M5 row допустимо и является диагностической provenance-информацией. Остальные object-level reject flags не меняются. Если требуется намеренно вернуть строгий series-level veto, его можно явно запросить через
\begin{lstlisting}
--extra-reject-flags SATURATED_LIKE
\end{lstlisting}
Это изменение относится только к default M5 row-selection; общий \texttt{DEFAULT\_REJECT\_FLAGS}, Stage~2 Object-QC и M1--M4 reference selection не меняются.

\subsection{Неопределённости M5}

В CSV сохраняются раздельно:

\begin{description}
  \item[\texttt{mag\_err\_random}] random component Stage~2 measurement;
  \item[\texttt{mag\_err\_science\_systematic}] Stage~2 systematic component, если он существует;
  \item[\texttt{mag\_err\_series\_zp}] M4 uncertainty $\sigma_{z,s}$, общая и коррелированная для всей series.
\end{description}

Для удобного одночислового отображения экспортируется
\[
  \sigma_{\mathrm{effective},i}
  =\sqrt{\sigma_{\mathrm{random},i}^2+\sigma_{z,s}^2}.
\]
Это поле является convenience estimate, а не независимой ошибкой каждой точки: $\sigma_{z,s}$ нельзя усреднять как white noise при последующем анализе.

\subsection{Ограничение числа точек на графике}

Основной CSV никогда не уменьшается из-за plot policy. Опция \texttt{--plot-max-points} относится только к визуализации и по умолчанию равна 5000 points на filter. Начиная с 0.5.47.2 это жёсткий верхний предел: итоговый binned plot не может содержать больше указанного числа points на filter. Каждый непустой series получает как минимум один bin; если заданный limit меньше числа непустых series, программа останавливает построение plot вместо скрытого удаления целых series.

В режиме \texttt{--plot-reduction auto} решение raw/bin консервативно принимается по числу exact linked source rows до QC; если оно не превышает лимит, строятся accepted исходные точки. Если лимит превышен, бюджет bins распределяется между series методом наибольших остатков с учётом их числа measurements, после чего каждая series бинингуется отдельно. Для series $s$ с выделенным числом bins $P_s$ автоматическая ширина приблизительно равна
\[
  \Delta t_s = \frac{t_{\max,s}-t_{\min,s}}{P_s},
\]
но не меньше одной секунды. Точка точно на $t_{\max,s}$ принудительно относится к последнему выделенному bin, поэтому inclusive endpoint не создаёт $(P_s+1)$-й bin. Можно задать физическую ширину через \texttt{--plot-bin-minutes}; в 0.5.47.2 она трактуется как минимальная ширина и при необходимости автоматически увеличивается, чтобы сохранить общий hard cap.

Для каждого непустого bin $B$ plot-layer вычисляет arithmetic mean
\[
  \bar m_B=\frac{1}{n_B}\sum_{i\in B}m_i,
\]
а random error для отображения выбирает консервативно:
\[
 \sigma_{B,\mathrm{rand}}=
 \max\left(
   \frac{s_B}{\sqrt{n_B}},
   \frac{\sqrt{\sum_{i\in B}\sigma_{\mathrm{random},i}^2}}{n_B}
 \right).
\]
Дополнительно сохраняются \texttt{mag\_min} и \texttt{mag\_max}; график показывает min--max envelope, поэтому сильные excursions не исчезают полностью при усреднении. M4 zero-point uncertainty на error bars не рисуется как независимая: она остаётся отдельным correlated series-level полем.

\texttt{--plot-reduction bin} принудительно включает биннинг, \texttt{--plot-reduction none} или \texttt{--no-plots} отключают визуализацию. Формат выбирается через \texttt{--plot-format png|pdf|svg}.

\subsection{Разрыв и сжатие оси времени}

Для multi-night overview физически пустые промежутки между сериями могут быть намного длиннее самих наблюдений. Опция
\begin{lstlisting}
--plot-time-axis auto|continuous|compressed
\end{lstlisting}
управляет только отображением. Значение \texttt{auto} используется по умолчанию: временные интервалы внутри каждого наблюдательного блока сохраняются без изменения, а большие gaps между неперекрывающимися series сокращаются до небольшого визуального зазора. На оси такие места отмечаются \texttt{//}. Подписи показывают исходный MJD, а при небольшом числе series также их идентификаторы.

\texttt{continuous} оставляет обычную линейную ось MJD. \texttt{compressed} принудительно предпочитает сокращение межсерийных gaps, когда это возможно. Это преобразование не изменяет \texttt{mjd} в основном CSV и \texttt{*\_plot.csv}; оно применяется только внутри Matplotlib renderer.

При числе series до 12 сохраняются отдельные визуальные стили и legend для series. Если series больше 12, renderer автоматически переходит в archive-overview: использует единый визуальный стиль внутри одного filter и убирает per-series legend, чтобы сотни цветов и подписей не занимали весь рисунок. Структура наблюдений при этом читается по MJD-подписям и отметкам разрывов времени.

\section{M8: resumable batch orchestration начиная с 0.5.48.0}

M8 не вводит новый photometric estimator и не меняет уравнения M1--M5. Это orchestration layer над уже реализованными scientific APIs. Он делит \texttt{processed-root} на leaf batches с FITS, сверяет их с БД и присваивает состояние \texttt{SKIP\_READY}, \texttt{RUN\_NEW}, \texttt{RESUME\_INCOMPLETE}, \texttt{SKIP\_TOO\_SMALL} или \texttt{STOP\_CONFLICT}. Начиная с 0.5.48.1 \texttt{SKIP\_TOO\_SMALL} выставляется до Stage~2, если ни одна series-like группа (profile/instrument/site/tube/camera/filter/object) не набирает \texttt{min\_frames\_per\_series=10}. READY данные при resume не перерабатываются. Leaf batch, содержащий FITS более чем одной календарной даты, намеренно блокируется как неоднозначный. Начиная с 0.5.54.3 readiness для P4-linked данных считается по физическим Observation, а не по обязательному прямому назначению каждого Frame product: alternate Frame не делает batch неполным, если та же Observation уже представлена READY science series данного field. При нескольких historical representatives M8 предпочитает READY representative; если READY representatives несколько, выбирается representative с наибольшим science-frame coverage. Observation без representative и legacy Frame с \texttt{observation\_id=NULL}, не назначенный напрямую, по-прежнему дают \texttt{RESUME\_INCOMPLETE}.

Stage~2 запускается frozen recipe \texttt{stage3\_scale\_v1}: \texttt{adaptive\_ellipse}, aperture-correction policy \texttt{evidence}, conventional ensemble с минимум 30 references/frame, candidate-excluded iteration до сходимости с growth policy \texttt{initial\_only}, photometric-frame-QC, PSF diagnostics и color-calibration attempt; spatial ensemble остаётся выключенным. Target provenance для архивного field-run отключается.

После новых Stage~2 products M8 последовательно требует:
\begin{enumerate}
  \item parity \texttt{science\_measurements} для новых/возобновлённых SCIENCE generations;
  \item \texttt{READY\_SHARED\_IDENTITY};
  \item M0 readiness для каждого явно заданного scope;
  \item connected M1--M3 graph со статусом \texttt{READY\_M1\_M3}/\texttt{OK};
  \item immutable M4 persist и \texttt{verify-interseries-calibration = OK};
  \item при заданных control objects --- M5 с ненулевым числом accepted points.
\end{enumerate}

Порог \texttt{min\_common\_objects} является частью explicit policy и никогда не изменяется orchestration автоматически. M8 v1 допускает ровно одну explicit policy на filter, чтобы выбор итоговой M4 generation оставался однозначным. Если graph disconnected, workflow останавливается и оставляет partial machine-readable report.

Итоговый M8 report bundle содержит \texttt{final\_report.json}, \texttt{report.html}, \texttt{batches.csv}, \texttt{frame\_qc.csv}, \texttt{series.csv}, \texttt{interseries\_edges.csv}, \texttt{interseries\_offsets.csv} и \texttt{control\_objects.csv}, а также полные M1--M3 diagnostics и M5 products в подкаталогах.

\section{Последующие этапы}

После стабилизации M8 следующий этап --- multi-series QC/variability metrics, затем support heterogeneous instrument units и cross-catalog identity.

Актуальная инженерная карта находится в \texttt{docs/ROADMAP.md}; исторический Stage~3 plan сохранён в \texttt{docs/PROJECT\_HISTORY.md}.

\section{Подготовка чистой Stage 2 базы перед M1}

Версия \texttt{0.5.41.1} добавляет промежуточный infrastructure milestone между M0 и M1. Его задача --- не менять фотометрический алгоритм, а обеспечить, чтобы Stage~3 сравнивал science products, пересчитанные одной современной версией Stage~2, а не смесь исторических benchmark generations.

После полного пересчёта следует проверить field-wide catalog identity:

\begin{lstlisting}
python -m grandgousier.cli audit-field-catalog-identity `
  --field-id 2
\end{lstlisting}

Команда группирует повторные series по \texttt{field\_id + instrument\_unit\_id + filter\_name} и проверяет catalog namespaces, реально связанные с кадрами через \texttt{frame\_catalog\_links}. Для прямого object-ID stitching первого Stage~3 MVP предпочтителен статус

\begin{lstlisting}
READY_SHARED_IDENTITY
\end{lstlisting}

и для каждой повторной группы

\begin{lstlisting}
READY_SHARED_NAMESPACE
\end{lstlisting}

Это означает, что серии одной фотометрической системы используют один и тот же catalog namespace. При текущей схеме \texttt{object\_uid} строится из catalog namespace/source identity, поэтому такой shared namespace обеспечивает стабильное повторное использование одного \texttt{object\_id} при новом forced-photometry проходе.

Статус \texttt{NEEDS\_RECONCILIATION} означает, что две series одной системы используют несовпадающие catalog namespaces. В таком состоянии M1/M2 запускать не следует: сначала нужна унификация catalog identity или будущий cross-catalog reconciliation.


\subsection{Стабильность catalog namespace при clean rebuild}

Stage~3 использует устойчивую object identity между series. Поэтому clean rebuild Stage~2 не должен без необходимости заново создавать уже существующий auto-catalog. Версия \texttt{0.5.41.2} повторно использует заполненный каталог с точным ожидаемым auto-generated именем до попытки Gaia query. Это сохраняет существующий namespace и исключает сетевую зависимость для уже известных полей. После rebuild всё равно необходимо выполнить \texttt{audit-field-catalog-identity}: reuse одного старого namespace не решает автоматически cross-catalog identity для series, которые исторически были построены в разных namespaces.

\subsection{Глобальный retirement старого SCIENCE V1 (0.6.4.5)}
После точного retirement отдельных полей команды обслуживания поддерживают явную область
\texttt{--all-materialized}. Она выбирает только неизменяемые SCIENCE-продукты, для которых
в базе ещё физически существуют строки \texttt{science\_measurements}. Флаг взаимоисключающий
с \texttt{--field-id} и явными \texttt{--calibration-version-id}; без \texttt{--execute}
никакого удаления не происходит. Типичная последовательность: заполнить compact cache,
построить и проверить exact flag overrides, выполнить полный retirement dry-run и только затем
повторить retirement с \texttt{--execute}. Уже retired продукты автоматически не попадают в
глобальную область, так как V1-строк у них больше нет.
